\documentclass[../DoAn.tex]{subfiles}
\begin{document}
\chapter{Khảo sát và phân tích yêu cầu}

Chương này trình bày quá trình khảo sát hiện trạng và phân tích yêu cầu cho hệ thống TheWeekend. Nội dung chương tập trung vào việc xác định bối cảnh nghiệp vụ, các nhóm tác nhân sử dụng hệ thống, các nhóm chức năng chính, các quy trình nghiệp vụ tiêu biểu và các yêu cầu phi chức năng ở mức định tính. Do chưa có dữ liệu khảo sát định lượng chính thức, chương này không đưa ra số liệu thống kê, kết quả khảo sát người dùng hoặc kết quả kiểm thử.

\section{Khảo sát hiện trạng}

Việc lựa chọn địa điểm vui chơi cuối tuần cho trẻ em và gia đình có trẻ em là một bài toán cần tổng hợp nhiều loại thông tin khác nhau. Khi xem xét một địa điểm, người dùng thường cần quan tâm đến vị trí, độ tuổi phù hợp, giá vé, giờ mở cửa, tiện ích, hình ảnh và đánh giá từ những người dùng khác. Nếu các thông tin này được trình bày rời rạc, người dùng phải tự tìm kiếm, đối chiếu và đánh giá trước khi đưa ra quyết định.

Trong phạm vi đồ án, quá trình khảo sát được thực hiện theo hướng phân tích hiện trạng nghiệp vụ và nhu cầu chức năng của hệ thống, chưa có dữ liệu khảo sát định lượng chính thức từ người dùng. Do đó, báo cáo không đưa ra các số liệu về số lượng người khảo sát, tỷ lệ lựa chọn hoặc mức độ hài lòng. Nội dung phân tích tập trung vào nhu cầu tổ chức dữ liệu địa điểm một cách tập trung, hỗ trợ người dùng tra cứu, xem chi tiết, lưu địa điểm yêu thích, đánh giá và nhận gợi ý phù hợp hơn với bối cảnh sử dụng.

Bên cạnh nhu cầu tra cứu thông tin, một số địa điểm có thể phát sinh nghiệp vụ đặt vé. Khi đó, hệ thống cần hỗ trợ quy trình từ chọn vé, tạo đơn đặt vé (booking), thanh toán, phát hành vé QR đến kiểm tra vé tại điểm đến. Việc số hóa quy trình này giúp dữ liệu đặt vé, thanh toán và trạng thái sử dụng vé được quản lý thống nhất trong hệ thống. Đây là cơ sở để xác định các nhóm chức năng chính của TheWeekend trong các phần tiếp theo.

\subsection{Phân tích so sánh các nền tảng tương tự}

Để xác định rõ hơn vị trí của TheWeekend, phần này thực hiện phân tích so sánh với một số nền tảng có chức năng liên quan đến tra cứu địa điểm, đặt vé hoặc phục vụ nhu cầu vui chơi. Đây là phân tích đặc điểm chức năng của các nền tảng tương tự dựa trên đặc điểm phổ biến đã biết của từng nền tảng, không phải kết quả khảo sát người dùng hay đo đạc định lượng. Các nền tảng được lựa chọn gồm: Google Maps (tra cứu địa điểm và bản đồ), Foody (tra cứu địa điểm ăn uống và đánh giá), Traveloka và Klook (đặt vé tham quan, vui chơi và du lịch), một nhóm website khu vui chơi cho trẻ em điển hình, và chính TheWeekend. Cột ``Web vui chơi trẻ em'' đại diện cho đặc điểm chung của các website giới thiệu khu vui chơi trẻ em hiện có, không trỏ tới một sản phẩm cụ thể.

Các tiêu chí so sánh được lựa chọn bám sát nghiệp vụ trọng tâm của TheWeekend, gồm: tra cứu địa điểm, lọc theo vị trí, lọc theo độ tuổi trẻ em, phân loại và gắn thẻ địa điểm, đánh giá của người dùng, gợi ý địa điểm theo ngữ cảnh, chatbot hỏi đáp dựa trên dữ liệu địa điểm, đặt vé và thanh toán, vé QR và check-in, quản trị địa điểm và vé. Mỗi tiêu chí được đánh giá ở ba mức: ``Có'' nếu nền tảng hỗ trợ rõ ràng, ``Một phần'' nếu hỗ trợ ở mức hạn chế hoặc gián tiếp, ``Không'' nếu không phải trọng tâm của nền tảng. Để dễ đọc, kết quả phân tích được tách thành hai bảng: Bảng~\ref{tab:platform-comparison-a} cho nhóm tiêu chí tra cứu và đánh giá, Bảng~\ref{tab:platform-comparison-b} cho nhóm tiêu chí gợi ý, đặt vé và quản trị. Trong cả hai bảng, cột ``Web trẻ em'' là viết tắt của nhóm website khu vui chơi trẻ em đã nêu ở trên.

\begin{table}[H]
\centering
\caption{So sánh nhóm tiêu chí tra cứu và đánh giá}
\label{tab:platform-comparison-a}
\scriptsize
\renewcommand{\arraystretch}{1.25}
\setlength{\tabcolsep}{3pt}
\begin{tabularx}{\textwidth}{|L{3.2cm}|*{6}{>{\centering\arraybackslash}X|}}
\hline
\textbf{Tiêu chí} & \textbf{Google Maps} & \textbf{Foody} & \textbf{Traveloka} & \textbf{Klook} & \textbf{Web trẻ em} & \textbf{TheWeekend} \\
\hline
Tra cứu địa điểm & Có & Có & Một phần & Có & Một phần & Có \\
\hline
Lọc theo vị trí & Có & Có & Có & Có & Một phần & Có \\
\hline
Lọc theo độ tuổi trẻ em & Không & Không & Không & Một phần & Một phần & Có \\
\hline
Phân loại và gắn thẻ & Có & Có & Có & Có & Một phần & Có \\
\hline
Đánh giá của người dùng & Có & Có & Có & Có & Một phần & Có \\
\hline
\end{tabularx}
\end{table}

\begin{table}[H]
\centering
\caption{So sánh nhóm tiêu chí gợi ý, đặt vé và quản trị}
\label{tab:platform-comparison-b}
\scriptsize
\renewcommand{\arraystretch}{1.25}
\setlength{\tabcolsep}{3pt}
\begin{tabularx}{\textwidth}{|L{3.2cm}|*{6}{>{\centering\arraybackslash}X|}}
\hline
\textbf{Tiêu chí} & \textbf{Google Maps} & \textbf{Foody} & \textbf{Traveloka} & \textbf{Klook} & \textbf{Web trẻ em} & \textbf{TheWeekend} \\
\hline
Gợi ý theo ngữ cảnh & Một phần & Một phần & Một phần & Một phần & Không & Có \\
\hline
Chatbot theo dữ liệu địa điểm & Không & Không & Một phần & Một phần & Không & Có \\
\hline
Đặt vé và thanh toán & Không & Một phần & Có & Có & Một phần & Có \\
\hline
Vé QR và check-in & Không & Không & Có & Có & Một phần & Có \\
\hline
Quản trị địa điểm và vé & Một phần & Có & Có & Có & Một phần & Có \\
\hline
\end{tabularx}
\end{table}

Từ Bảng~\ref{tab:platform-comparison-a} và Bảng~\ref{tab:platform-comparison-b}, có thể thấy các nền tảng phổ biến đều mạnh ở một số nhóm chức năng nhất định nhưng ít nền tảng tập trung đồng thời vào nhu cầu chọn địa điểm vui chơi cho trẻ em. Google Maps và Foody mạnh về tra cứu địa điểm và đánh giá nhưng không hướng tới lọc theo độ tuổi trẻ em hay đặt vé. Traveloka và Klook hỗ trợ tốt đặt vé, vé QR và check-in nhưng phạm vi thiên về du lịch và sự kiện nói chung, chưa chuyên biệt cho tiêu chí phù hợp với trẻ em. Các website khu vui chơi trẻ em thường chỉ giới thiệu thông tin ở mức cơ bản, ít có gợi ý theo ngữ cảnh, chatbot hoặc quy trình vé điện tử hoàn chỉnh. TheWeekend hướng tới kết hợp các nhóm chức năng này trong cùng một hệ thống: tra cứu và lọc địa điểm theo độ tuổi, gợi ý dựa trên dữ liệu nội bộ, chatbot hỏi đáp, đặt vé, vé QR và check-in. Phân tích này cho thấy khoảng trống chức năng mà TheWeekend nhắm tới, đồng thời lưu ý rằng đánh giá ở đây mang tính so sánh đặc điểm chức năng, chưa bao hàm so sánh về quy mô dữ liệu, hiệu năng hay trải nghiệm thực tế của các nền tảng.

\section{Tổng quan chức năng}

Hệ thống TheWeekend phục vụ bốn nhóm tác nhân chính: khách chưa đăng nhập, người dùng đã đăng nhập, quản trị viên và nhân viên. Mỗi nhóm tác nhân có phạm vi sử dụng khác nhau tùy theo vai trò trong hệ thống. Bảng~\ref{tab:actors} trình bày tổng quan các actor và nhóm chức năng chính.

\begin{table}[H]
\centering
\caption{Actor và nhóm chức năng chính trong hệ thống TheWeekend}
\label{tab:actors}
\small
\renewcommand{\arraystretch}{1.15}
\setlength{\tabcolsep}{4pt}
\begin{tabularx}{\textwidth}{|L{2.6cm}|Y|Y|L{3.0cm}|}
\hline
\textbf{Actor} & \textbf{Mô tả} & \textbf{Nhóm chức năng được sử dụng} & \textbf{Chức năng liên quan/Ghi chú} \\
\hline
Khách chưa đăng nhập & Người chưa có phiên đăng nhập trong hệ thống & Đăng ký, đăng nhập, đăng nhập Google, quên mật khẩu, đặt lại mật khẩu, xem một số thông tin công khai & Tài khoản và tra cứu công khai \\
\hline
Người dùng đã đăng nhập & Người dùng có tài khoản và vai trò người dùng thông thường & Quản lý hồ sơ, tìm kiếm, xem chi tiết địa điểm, nhận gợi ý, chatbot, yêu thích, đánh giá, gửi ảnh đánh giá, đặt vé, thanh toán, xem vé & Chức năng người dùng cuối \\
\hline
Quản trị viên & Người dùng có quyền quản trị hệ thống & Quản lý người dùng, địa điểm, đánh giá, ảnh đánh giá, vé, thanh toán và check-in & Quản trị và vận hành \\
\hline
Nhân viên & Người dùng được cấp quyền phục vụ nghiệp vụ tại điểm đến & Quét QR, xác minh vé, check-in vé, xem một số nghiệp vụ vé được cấp quyền & Check-in vé \\
\hline
\end{tabularx}
\end{table}

Từ Bảng~\ref{tab:actors}, có thể thấy hệ thống được phân tách theo vai trò rõ ràng. Người dùng thông thường tập trung vào các chức năng tra cứu, gợi ý, đánh giá và đặt vé. Quản trị viên đảm nhiệm việc quản lý dữ liệu và theo dõi nghiệp vụ. Nhân viên tham gia vào quy trình vận hành tại điểm đến thông qua chức năng quét QR và check-in vé.

\subsection{Biểu đồ use case tổng quát}

Biểu đồ use case tổng quát trong Hình~\ref{fig:usecase-overview} mô tả các nhóm chức năng chính của hệ thống TheWeekend. Để biểu đồ dễ đọc, các use case được gom theo nhóm nghiệp vụ thay vì liệt kê toàn bộ chức năng chi tiết.

\begin{figure}[H]
\centering
\includegraphics[width=\textwidth,height=0.85\textheight,keepaspectratio]{figures/usecase_overview.png}
\caption{Biểu đồ use case tổng quát hệ thống TheWeekend}
\label{fig:usecase-overview}
\end{figure}

Trong Hình~\ref{fig:usecase-overview}, khách chưa đăng nhập chủ yếu tương tác với nhóm chức năng tài khoản và tra cứu thông tin công khai. Người dùng đã đăng nhập sử dụng các chức năng chính của hệ thống như tra cứu địa điểm, nhận gợi ý, hỏi đáp, đánh giá và đặt vé. Quản trị viên quản lý dữ liệu hệ thống và nghiệp vụ vé. Nhân viên tham gia vào quy trình xác minh và check-in vé QR.

\subsection{Biểu đồ use case phân rã}

Để mô tả rõ hơn các nghiệp vụ trọng tâm, báo cáo lựa chọn ba biểu đồ use case phân rã đưa vào Chương 2, gồm nhóm địa điểm và gợi ý, nhóm đặt vé và thanh toán, nhóm quản trị và check-in. Nhóm đánh giá có thể được mô tả trong phụ lục nếu cần trình bày chi tiết hơn, nhằm tránh làm phần phân tích yêu cầu quá dài.

Biểu đồ phân rã nhóm địa điểm và gợi ý trong Hình~\ref{fig:usecase-place-suggestion} mô tả các chức năng giúp người dùng tra cứu, xem thông tin và nhận đề xuất địa điểm. Actor chính gồm khách chưa đăng nhập và người dùng đã đăng nhập. Người dùng đã đăng nhập có thêm các thao tác cá nhân hóa như cập nhật vị trí và lưu địa điểm yêu thích.

\begin{figure}[H]
\centering
\includegraphics[width=\textwidth,height=0.85\textheight,keepaspectratio]{figures/usecase_place_suggestion.png}
\caption{Biểu đồ use case phân rã nhóm địa điểm và gợi ý}
\label{fig:usecase-place-suggestion}
\end{figure}

Biểu đồ phân rã nhóm đặt vé và thanh toán trong Hình~\ref{fig:usecase-ticket-payment} mô tả chuỗi chức năng từ chọn vé đến theo dõi trạng thái thanh toán và nhận vé QR. Actor chính là người dùng đã đăng nhập; quản trị viên tham gia theo dõi và quản lý nghiệp vụ vé.

\begin{figure}[H]
\centering
\includegraphics[width=\textwidth,height=0.85\textheight,keepaspectratio]{figures/usecase_ticket_payment.png}
\caption{Biểu đồ use case phân rã nhóm đặt vé và thanh toán}
\label{fig:usecase-ticket-payment}
\end{figure}

Biểu đồ phân rã nhóm quản trị và check-in trong Hình~\ref{fig:usecase-admin-checkin} mô tả các chức năng dành cho quản trị viên và nhân viên. Nhóm này bao gồm quản lý dữ liệu hệ thống và thao tác vận hành vé tại điểm đến.

\begin{figure}[H]
\centering
\includegraphics[width=\textwidth,height=0.85\textheight,keepaspectratio]{figures/usecase_admin_checkin.png}
\caption{Biểu đồ use case phân rã nhóm quản trị và check-in}
\label{fig:usecase-admin-checkin}
\end{figure}

\subsection{Quy trình nghiệp vụ}

Trong phạm vi Chương 2, báo cáo lựa chọn bốn quy trình nghiệp vụ tiêu biểu để mô tả bằng activity diagram: tìm kiếm và gợi ý địa điểm, đặt vé và thanh toán PayOS, quét QR và check-in vé, hỏi đáp chatbot dựa trên dữ liệu nội bộ. Các quy trình này thể hiện rõ các nghiệp vụ đặc trưng của hệ thống TheWeekend, trong khi phần phân tích kỹ thuật chi tiết sẽ được trình bày ở Chương 4.

Quy trình tìm kiếm và gợi ý địa điểm bắt đầu khi người dùng nhập từ khóa, chọn tiêu chí tìm kiếm hoặc mở trang gợi ý. Hệ thống tiếp nhận tiêu chí, truy vấn dữ liệu địa điểm, sử dụng thông tin danh mục, tag hoặc vị trí nếu có, sau đó trả về danh sách phù hợp hoặc thông báo không có dữ liệu. Hình~\ref{fig:activity-place-suggestion} mô tả tổng quan quy trình này.

\begin{figure}[H]
\centering
\IfFileExists{figures/activity_place_suggestion.png}{\includegraphics[width=\textwidth,height=0.85\textheight,keepaspectratio]{figures/activity_place_suggestion.png}}{\fbox{\parbox{0.9\textwidth}{\centering [Hình chưa sẵn sàng: \texttt{activity\_place\_suggestion.png}]}}}
\caption{Activity diagram quy trình tìm kiếm và gợi ý địa điểm}
\label{fig:activity-place-suggestion}
\end{figure}

Quy trình đặt vé bắt đầu khi người dùng chọn địa điểm, ngày đi và số lượng vé. Hệ thống kiểm tra dữ liệu đầu vào, tạo đơn đặt vé, tạo yêu cầu thanh toán qua PayOS và cập nhật trạng thái đơn sau khi thanh toán được xác nhận. Hình~\ref{fig:activity-payment} mô tả tổng quan quy trình này.

\begin{figure}[H]
\centering
\includegraphics[width=\textwidth,height=0.85\textheight,keepaspectratio]{figures/activity_payment.png}
\caption{Activity diagram quy trình đặt vé và thanh toán PayOS}
\label{fig:activity-payment}
\end{figure}

Quy trình check-in vé được thực hiện bởi nhân viên hoặc quản trị viên thông qua ứng dụng Android scanner. Ứng dụng quét mã QR, gửi dữ liệu vé tới backend để xác minh, sau đó cho phép ghi nhận check-in nếu vé hợp lệ. Hình~\ref{fig:activity-checkin} mô tả quy trình này.

\begin{figure}[H]
\centering
\includegraphics[width=\textwidth,height=0.85\textheight,keepaspectratio]{figures/activity_checkin.png}
\caption{Activity diagram quy trình quét QR và check-in vé}
\label{fig:activity-checkin}
\end{figure}

Quy trình hỏi đáp chatbot bắt đầu khi người dùng gửi câu hỏi. Hệ thống sử dụng dữ liệu nội bộ về địa điểm, đánh giá và thẻ để tạo ngữ cảnh trả lời, sau đó kết hợp Gemini API để sinh phản hồi cho người dùng. Hình~\ref{fig:activity-chatbot} mô tả quy trình nghiệp vụ ở mức tổng quan.

\begin{figure}[H]
\centering
\includegraphics[width=\textwidth,height=0.85\textheight,keepaspectratio]{figures/activity_chatbot.png}
\caption{Activity diagram quy trình hỏi đáp chatbot}
\label{fig:activity-chatbot}
\end{figure}

\section{Đặc tả chức năng}

Phần này đặc tả năm use case quan trọng nhất của hệ thống. Các use case được lựa chọn vì thể hiện trực tiếp nghiệp vụ cốt lõi của TheWeekend, bao gồm tra cứu địa điểm, gợi ý, chatbot, đặt vé và check-in vé. Các chức năng đánh giá địa điểm và gửi ảnh đánh giá đã được mô tả ở mức tổng quan trong biểu đồ use case; nếu cần trình bày chi tiết hơn, phần đặc tả có thể được đưa sang phụ lục để tránh làm Chương 2 quá dài. Các bảng đặc tả được trình bày ngắn gọn để làm rõ mục tiêu, điều kiện, luồng xử lý và kết quả của từng use case.

\begin{table}[H]
\centering
\caption{Đặc tả use case tìm kiếm và xem chi tiết địa điểm}
\label{tab:uc-search-detail}
\small
\renewcommand{\arraystretch}{1.15}
\setlength{\tabcolsep}{4pt}
\begin{tabularx}{\textwidth}{|L{3.4cm}|Y|}
\hline
\textbf{Thuộc tính} & \textbf{Nội dung} \\
\hline
Tên use case & Tìm kiếm và xem chi tiết địa điểm \\
\hline
Actor chính & Khách, người dùng đã đăng nhập \\
\hline
Mục tiêu & Tìm địa điểm phù hợp và xem thông tin chi tiết \\
\hline
Tiền điều kiện & Hệ thống có dữ liệu địa điểm \\
\hline
Luồng chính & Người dùng nhập từ khóa hoặc chọn tiêu chí tìm kiếm; hệ thống truy vấn danh sách địa điểm; người dùng chọn một địa điểm; hệ thống hiển thị thông tin chi tiết \\
\hline
Luồng thay thế / ngoại lệ & Không có kết quả phù hợp; địa điểm không tồn tại; lỗi lấy dữ liệu \\
\hline
Hậu điều kiện & Người dùng xem được thông tin địa điểm hoặc nhận thông báo phù hợp \\
\hline
Dữ liệu đầu vào & Từ khóa, vị trí, danh mục, thẻ, độ tuổi nếu có \\
\hline
Kết quả đầu ra & Danh sách địa điểm và thông tin chi tiết địa điểm \\
\hline
\end{tabularx}
\end{table}

\begin{table}[H]
\centering
\caption{Đặc tả use case nhận gợi ý địa điểm}
\label{tab:uc-suggestion}
\small
\renewcommand{\arraystretch}{1.15}
\setlength{\tabcolsep}{4pt}
\begin{tabularx}{\textwidth}{|L{3.4cm}|Y|}
\hline
\textbf{Thuộc tính} & \textbf{Nội dung} \\
\hline
Tên use case & Nhận gợi ý địa điểm \\
\hline
Actor chính & Người dùng đã đăng nhập \\
\hline
Mục tiêu & Nhận danh sách địa điểm phù hợp với ngữ cảnh sử dụng \\
\hline
Tiền điều kiện & Người dùng đã đăng nhập; hệ thống có dữ liệu địa điểm \\
\hline
Luồng chính & Người dùng truy cập trang gợi ý; hệ thống lấy thông tin hồ sơ hoặc vị trí nếu có; hệ thống xử lý dữ liệu địa điểm; hệ thống hiển thị danh sách gợi ý \\
\hline
Luồng thay thế / ngoại lệ & Người dùng chưa có vị trí; không có dữ liệu phù hợp; lỗi tải dữ liệu \\
\hline
Hậu điều kiện & Người dùng nhận được danh sách địa điểm gợi ý hoặc thông báo không có dữ liệu \\
\hline
Dữ liệu đầu vào & Hồ sơ người dùng, vị trí, danh mục, dữ liệu địa điểm \\
\hline
Kết quả đầu ra & Danh sách địa điểm gợi ý \\
\hline
\end{tabularx}
\end{table}

\begin{table}[H]
\centering
\caption{Đặc tả use case hỏi đáp với chatbot}
\label{tab:uc-chatbot}
\small
\renewcommand{\arraystretch}{1.15}
\setlength{\tabcolsep}{4pt}
\begin{tabularx}{\textwidth}{|L{3.4cm}|Y|}
\hline
\textbf{Thuộc tính} & \textbf{Nội dung} \\
\hline
Tên use case & Hỏi đáp với chatbot \\
\hline
Actor chính & Người dùng \\
\hline
Mục tiêu & Nhận câu trả lời về địa điểm dựa trên dữ liệu nội bộ \\
\hline
Tiền điều kiện & API chatbot hoạt động; hệ thống có dữ liệu địa điểm, đánh giá hoặc thẻ \\
\hline
Luồng chính & Người dùng nhập câu hỏi; hệ thống truy hồi dữ liệu liên quan từ MongoDB; hệ thống tạo ngữ cảnh; hệ thống gọi Gemini API; chatbot trả lời và gợi ý địa điểm liên quan nếu có \\
\hline
Luồng thay thế / ngoại lệ & Câu hỏi trống; không có dữ liệu phù hợp; lỗi khi gọi Gemini API; hệ thống trả phản hồi không có dữ liệu hoặc phản hồi dự phòng \\
\hline
Hậu điều kiện & Phiên hội thoại được cập nhật \\
\hline
Dữ liệu đầu vào & Câu hỏi, mã hội thoại, vị trí người dùng nếu có \\
\hline
Kết quả đầu ra & Câu trả lời, danh sách địa điểm liên quan, gợi ý câu hỏi tiếp theo \\
\hline
\end{tabularx}
\end{table}

\begin{table}[H]
\centering
\caption{Đặc tả use case đặt vé và thanh toán}
\label{tab:uc-payment}
\small
\renewcommand{\arraystretch}{1.15}
\setlength{\tabcolsep}{4pt}
\begin{tabularx}{\textwidth}{|L{3.4cm}|Y|}
\hline
\textbf{Thuộc tính} & \textbf{Nội dung} \\
\hline
Tên use case & Đặt vé và thanh toán \\
\hline
Actor chính & Người dùng đã đăng nhập \\
\hline
Mục tiêu & Tạo đơn đặt vé, thanh toán qua PayOS và nhận vé QR \\
\hline
Tiền điều kiện & Người dùng đã đăng nhập; địa điểm có loại vé hợp lệ; PayOS được cấu hình \\
\hline
Luồng chính & Người dùng chọn địa điểm, ngày đi và số lượng vé; hệ thống tạo đơn đặt vé; hệ thống tạo yêu cầu thanh toán PayOS; người dùng thanh toán; hệ thống xác nhận trạng thái; hệ thống sinh vé QR \\
\hline
Luồng thay thế / ngoại lệ & Ngày đi không hợp lệ; địa điểm chưa có vé; thanh toán hết hạn; số tiền không khớp; webhook không hợp lệ \\
\hline
Hậu điều kiện & Đơn đặt vé được cập nhật trạng thái; vé QR được tạo khi thanh toán thành công \\
\hline
Dữ liệu đầu vào & Địa điểm, ngày đi, loại vé, số lượng vé, ghi chú \\
\hline
Kết quả đầu ra & Đơn đặt vé, thông tin thanh toán, vé QR sau khi thanh toán thành công \\
\hline
\end{tabularx}
\end{table}

\begin{table}[H]
\centering
\caption{Đặc tả use case quét QR và check-in vé}
\label{tab:uc-checkin}
\small
\renewcommand{\arraystretch}{1.15}
\setlength{\tabcolsep}{4pt}
\begin{tabularx}{\textwidth}{|L{3.4cm}|Y|}
\hline
\textbf{Thuộc tính} & \textbf{Nội dung} \\
\hline
Tên use case & Quét QR và check-in vé \\
\hline
Actor chính & Nhân viên, quản trị viên \\
\hline
Mục tiêu & Xác minh vé và ghi nhận lượt sử dụng vé tại điểm đến \\
\hline
Tiền điều kiện & Nhân viên hoặc quản trị viên đã đăng nhập; vé QR đã được phát hành \\
\hline
Luồng chính & Actor mở ứng dụng Android scanner; quét mã QR; ứng dụng gửi mã vé tới backend; hệ thống xác minh vé; actor xác nhận check-in; hệ thống cập nhật trạng thái vé \\
\hline
Luồng thay thế / ngoại lệ & QR không hợp lệ; vé không tồn tại; vé đã sử dụng; vé không ở trạng thái hợp lệ; lỗi kết nối \\
\hline
Hậu điều kiện & Vé được ghi nhận là đã sử dụng nếu hợp lệ \\
\hline
Dữ liệu đầu vào & QR payload hoặc mã vé \\
\hline
Kết quả đầu ra & Trạng thái xác minh vé và kết quả check-in \\
\hline
\end{tabularx}
\end{table}

\section{Yêu cầu phi chức năng}

Bên cạnh các yêu cầu chức năng, hệ thống cần đáp ứng một số yêu cầu phi chức năng để bảo đảm khả năng vận hành, bảo mật, bảo trì và tích hợp. Các yêu cầu được trình bày ở mức mục tiêu thiết kế và ràng buộc nghiệp vụ. Do chưa có kết quả kiểm thử định lượng chính thức, báo cáo không khẳng định các chỉ tiêu như thời gian phản hồi, số lượng người dùng đồng thời hoặc mức tải tối đa là kết quả đo đạc; các nội dung liên quan đến hiệu năng được phát biểu dưới dạng yêu cầu thiết kế trong phạm vi dữ liệu của đồ án. Bảng~\ref{tab:nonfunctional} trình bày các nhóm yêu cầu phi chức năng cùng mức độ đáp ứng tương ứng.

\begingroup
\small
\renewcommand{\arraystretch}{1.2}
\setlength{\tabcolsep}{4pt}
\begin{longtable}{|L{3.0cm}|L{8.4cm}|L{2.8cm}|}
\caption{Yêu cầu phi chức năng của hệ thống}\label{tab:nonfunctional}\\
\hline
\textbf{Nhóm yêu cầu} & \textbf{Nội dung yêu cầu (mục tiêu thiết kế)} & \textbf{Mức độ đáp ứng} \\
\hline
\endfirsthead
\hline
\textbf{Nhóm yêu cầu} & \textbf{Nội dung yêu cầu (mục tiêu thiết kế)} & \textbf{Mức độ đáp ứng} \\
\hline
\endhead
Thời gian phản hồi tìm kiếm và gợi ý & Hệ thống hướng tới trả về kết quả tìm kiếm và danh sách gợi ý trong thời gian ngắn, đủ để duy trì trải nghiệm tra cứu liên tục trong phạm vi dữ liệu của đồ án. Đây là mục tiêu thiết kế, không phải kết quả đo hiệu năng. & Mục tiêu thiết kế; chưa đo định lượng \\
\hline
Bảo mật tài khoản & Mật khẩu người dùng được băm bằng bcrypt và không lưu mật khẩu gốc; phiên đăng nhập sử dụng JWT để xác thực giữa client và backend. & Đã triển khai \\
\hline
Kiểm soát quyền truy cập & Hệ thống phân quyền theo vai trò người dùng, quản trị viên và nhân viên; áp dụng route guard ở phía giao diện và middleware kiểm tra quyền ở backend cho các API cần bảo vệ. & Đã triển khai \\
\hline
An toàn thanh toán & Hệ thống không tin tưởng trạng thái thanh toán do frontend gửi lên; trạng thái đã thanh toán chỉ được cập nhật sau khi xác minh dữ liệu từ PayOS (kiểm tra chữ ký webhook và đối chiếu số tiền) ở phía backend. & Đã triển khai \\
\hline
Độ tin cậy phát hành vé QR & Vé QR chỉ được phát hành sau khi giao dịch thanh toán được xác nhận thành công, nhằm tránh phát hành vé cho đơn chưa thanh toán. & Đã triển khai \\
\hline
Chống dùng lại vé QR & Vé đã được ghi nhận sử dụng (used) sẽ không được chấp nhận check-in lần nữa; thao tác quét lại vé đã sử dụng bị từ chối. & Đã triển khai \\
\hline
Khả năng mở rộng dữ liệu địa điểm & Dữ liệu địa điểm được tổ chức cùng danh mục (category) và thẻ (tag) để hỗ trợ phân loại và tra cứu; việc bổ sung chỉ mục (index) phục vụ tìm kiếm theo tên, địa chỉ, tag và danh mục được đề xuất như một hướng tối ưu. & Một phần; index tìm kiếm là đề xuất tối ưu \\
\hline
Sao lưu dữ liệu & Hệ thống cần có cơ chế sao lưu dữ liệu định kỳ khi sử dụng thực tế. Trong phạm vi đồ án, đây được nêu ở mức yêu cầu, chưa khẳng định đã có quy trình sao lưu định kỳ. & Yêu cầu \\
\hline
Tính dễ sử dụng & Giao diện cần hỗ trợ rõ các luồng chính như tìm kiếm, xem chi tiết, gợi ý, đặt vé, xem vé và quản trị; phản hồi thao tác được hiển thị sau các bước cần xác nhận. & Đã triển khai ở mức giao diện \\
\hline
Tính dễ bảo trì & Mã nguồn được tổ chức theo module (route, model, service, middleware ở backend; trang và component ở frontend) để thuận tiện mở rộng và sửa đổi. & Đã triển khai \\
\hline
Khả năng tích hợp dịch vụ ngoài & Hệ thống tích hợp các dịch vụ ngoài phục vụ bản đồ (Goong API), chatbot (Gemini API), thanh toán (PayOS) và email thông qua tầng service của backend. & Đã triển khai \\
\hline
Khả năng ghi nhận trạng thái giao dịch và vé & Hệ thống lưu vết trạng thái thanh toán, vé và check-in (kèm log sự kiện thanh toán) để phục vụ theo dõi nghiệp vụ. & Đã triển khai \\
\hline
\end{longtable}
\endgroup

Bảng~\ref{tab:nonfunctional} cho thấy phần lớn các yêu cầu phi chức năng liên quan đến bảo mật, phân quyền và toàn vẹn nghiệp vụ vé/thanh toán đã được hiện thực hóa ở mức chức năng, trong khi các yêu cầu liên quan đến hiệu năng định lượng, chỉ mục tìm kiếm và sao lưu dữ liệu được phát biểu ở mức mục tiêu thiết kế. Các yêu cầu định lượng như thời gian phản hồi đo đạc, tải đồng thời và độ ổn định dài hạn cần được bổ sung sau khi có kế hoạch kiểm thử và kết quả đo đạc chính thức.

Chương này đã trình bày hiện trạng nghiệp vụ, các nhóm actor, tổng quan chức năng, các use case chính, một số quy trình nghiệp vụ quan trọng và yêu cầu phi chức năng của hệ thống TheWeekend. Các nội dung trên là cơ sở để lựa chọn công nghệ ở Chương 3 và trình bày phân tích thiết kế, triển khai hệ thống ở Chương 4.
\end{document}




